% ============================================================
% Breadcrumb navigation and footer chrome
%
% Shared metadata-driven API:
%   \lrasetup{edition=<digital|print>, paper=<letter|a4|sixbynine>}
%   \lrameta{series={...}, volume={...}, book={...}, chapter={...},
%            topic={...}, current=<series|volume|book|chapter|topic>}
%   \LraBreadcrumb
%   \LraFooterBreadcrumb
%
% Compatibility API retained during migration:
%   \breadcrumb{<chapter subject>}{<prior display>}{<current display>}{<next display>}
% ============================================================

\makeatletter
\RequirePackage{fancyhdr}
\RequirePackage{xparse}

% ------------------------------------------------------------
% Configuration and metadata store
% ------------------------------------------------------------
\newif\iflra@icons
\lra@iconstrue
\newif\iflra@body
\lra@bodyfalse

\def\lra@edition{digital}
\def\lra@paper{letter}
\def\lra@threshold{360pt}
\def\lra@series{}
\def\lra@volume{}
\def\lra@book{}
\def\lra@chapter{}
\def\lra@section{}
\def\lra@topic{}
\def\lra@current{topic}

\pgfkeys{
  /lra/.is family, /lra,
  edition/.store in = \lra@edition,
  paper/.store in = \lra@paper,
  threshold/.store in = \lra@threshold,
  series/.store in = \lra@series,
  volume/.store in = \lra@volume,
  book/.store in = \lra@book,
  chapter/.store in = \lra@chapter,
  section/.store in = \lra@section,
  topic/.store in = \lra@topic,
  current/.store in = \lra@current,
  icons/.is choice,
  icons/true/.code = {\lra@iconstrue},
  icons/false/.code = {\lra@iconsfalse},
}

\newcommand{\lra@applypaper}{%
  \edef\lra@pcur{\lra@paper}%
  \def\lra@tmp{a4}%
  \ifx\lra@pcur\lra@tmp
    \geometry{a4paper,margin=2.2cm}%
  \fi
  \def\lra@tmp{letter}%
  \ifx\lra@pcur\lra@tmp
    \geometry{letterpaper,margin=1in}%
  \fi
  \def\lra@tmp{sixbynine}%
  \ifx\lra@pcur\lra@tmp
    \geometry{paperwidth=6in,paperheight=9in,margin=0.75in}%
  \fi
}

\newcommand{\lrasetup}[1]{%
  \pgfkeys{/lra,#1}%
  \iflra@body\else
    \lra@applypaper
  \fi
}
\newcommand{\lrameta}[1]{\pgfkeys{/lra,#1}}

% Read build-provided defaults. The build helper passes LRA_EDITION and
% LRA_PAPER; LRA_PRINT_EDITION remains accepted for legacy build commands.
\ifdefined\directlua
  \directlua{
    local edition = os.getenv("LRA_EDITION")
    local paper = os.getenv("LRA_PAPER")
    local legacy_print = os.getenv("LRA_PRINT_EDITION")
    if legacy_print and legacy_print ~= "" and legacy_print ~= "0"
       and legacy_print:lower() ~= "false" and legacy_print:lower() ~= "no" then
      edition = "print"
    end
    if edition == "digital" or edition == "print" then
      tex.sprint("\\lrasetup{edition=" .. edition .. "}")
    end
    if paper == "letter" or paper == "a4" or paper == "sixbynine" then
      tex.sprint("\\lrasetup{paper=" .. paper .. "}")
    end
  }
\fi
\AtBeginDocument{\lra@applypaper\global\lra@bodytrue}

% ------------------------------------------------------------
% Color system: identity hue plus state treatment
% ------------------------------------------------------------
\definecolor{lraSeriesFull}{HTML}{0A3B38}
\definecolor{lraSeriesTint}{HTML}{DCE8E6}
\definecolor{lraVolumeFull}{HTML}{176B66}
\definecolor{lraVolumeTint}{HTML}{DBEAE8}
\definecolor{lraBookFull}{HTML}{8A5A12}
\definecolor{lraBookTint}{HTML}{F1E6CF}
\definecolor{lraChapterFull}{HTML}{5A4E44}
\definecolor{lraChapterTint}{HTML}{E7E2DC}
\definecolor{lraSectionFull}{HTML}{0F5F73}
\definecolor{lraSectionTint}{HTML}{D8EAF0}
\definecolor{lraTopicFull}{HTML}{6E2A38}
\definecolor{lraTopicTint}{HTML}{EFD9DD}

% ------------------------------------------------------------
% Render-stable TikZ icons, taken from the breadcrumb package draft.
% ------------------------------------------------------------
\newcommand{\lra@iconbox}[1]{%
  \raisebox{-0.2em}{%
    \begin{tikzpicture}[x=1.5pt,y=1.5pt,line width=0.5pt,line cap=round,line join=round]%
      #1%
    \end{tikzpicture}%
  }%
}
\newcommand{\lraicon@series}{\lra@iconbox{%
  \draw (0,3)--(4,5)--(8,3)--(4,1)--cycle;
  \draw (0,4.5)--(4,6.5)--(8,4.5);
  \draw (0,6)--(4,8)--(8,6);}}
\newcommand{\lraicon@volume}{\lra@iconbox{%
  \draw (1,0) rectangle (8,8);\draw (2.6,0)--(2.6,8);}}
\newcommand{\lraicon@book}{\lra@iconbox{%
  \draw (4,1.6)--(4,7.4);
  \draw (4,7.4) .. controls (2.4,6.6) and (0.8,6.7) .. (0.4,6.9)--(0.4,1.5)
        .. controls (1.2,1.2) and (3,1.4) .. (4,2);
  \draw (4,7.4) .. controls (5.6,6.6) and (7.2,6.7) .. (7.6,6.9)--(7.6,1.5)
        .. controls (6.8,1.2) and (5,1.4) .. (4,2);}}
\newcommand{\lraicon@chapter}{\lra@iconbox{%
  \draw (0.5,6.8)--(7.5,6.8);\draw (0.5,4.7)--(5.3,4.7);
  \draw (0.5,2.6)--(7,2.6);\draw (0.5,0.7)--(3.8,0.7);}}
\newcommand{\lraicon@section}{\lra@iconbox{%
  \draw (0.7,6.8)--(7.3,6.8);\draw (0.7,4.8)--(7.3,4.8);
  \draw (0.7,2.8)--(5.8,2.8);\draw (0.7,0.8)--(4.2,0.8);}}
\newcommand{\lraicon@topic}{\lra@iconbox{%
  \draw (1,0)--(1,8)--(5.4,8)--(7,6.4)--(7,0)--cycle;
  \draw (5.4,8)--(5.4,6.4)--(7,6.4);
  \draw (2.2,5.4)--(5.8,5.4);\draw (2.2,3.8)--(5.8,3.8);\draw (2.2,2.2)--(4.6,2.2);}}

% ------------------------------------------------------------
% Metadata parsing helpers
% ------------------------------------------------------------
\newcommand{\lra@levelnum}[1]{%
  \edef\lra@tmp{#1}%
  \def\lra@N{6}%
  \def\lra@t{series}\ifx\lra@tmp\lra@t\def\lra@N{1}\fi
  \def\lra@t{volume}\ifx\lra@tmp\lra@t\def\lra@N{2}\fi
  \def\lra@t{book}\ifx\lra@tmp\lra@t\def\lra@N{3}\fi
  \def\lra@t{chapter}\ifx\lra@tmp\lra@t\def\lra@N{4}\fi
  \def\lra@t{section}\ifx\lra@tmp\lra@t\def\lra@N{5}\fi
  \def\lra@t{topic}\ifx\lra@tmp\lra@t\def\lra@N{6}\fi
}
\def\lra@splitcolon#1: #2\relax{\def\lra@eye{#1}\def\lra@ttl{#2}}
\newcommand{\lra@parse}[2]{%
  \def\lra@eye{#2}%
  \def\lra@ttl{#1}%
  \in@{: }{#1}%
  \ifin@
    \lra@splitcolon#1\relax
  \fi
}
\def\lra@titleget#1{%
  \in@{: }{#1}%
  \ifin@
    \lra@splitcolon#1\relax
    \lra@ttl
  \else
    #1%
  \fi
}
\newcommand{\lra@titleof}[1]{\expandafter\lra@titleget\expandafter{#1}}

% ------------------------------------------------------------
% Metadata-driven breadcrumb renderers
% ------------------------------------------------------------
\newcommand{\LraBreadcrumb}{%
  \ifdim\linewidth<\lra@threshold
    \LraCompactBreadcrumb
  \else
    \lra@chevron
  \fi
}

\newcommand{\lra@chevron}{%
  \lra@levelnum{\lra@current}%
  \begingroup
  \pgfmathsetmacro{\lra@d}{8}%
  \pgfmathsetmacro{\lra@W}{\linewidth/1pt}%
  \pgfmathsetmacro{\lra@w}{(\lra@W-\lra@d)/\lra@N}%
  \pgfmathsetmacro{\lra@h}{46}%
  \noindent\begin{tikzpicture}[x=1pt,y=1pt]
    \foreach \i in {1,...,\lra@N}{\lra@segment{\i}}
  \end{tikzpicture}%
  \endgroup
}

\newcommand{\lra@segment}[1]{%
  \pgfmathsetmacro{\xl}{(#1-1)*\lra@w}%
  \pgfmathsetmacro{\xr}{\xl+\lra@w}%
  \pgfmathsetmacro{\xrt}{\xr+\lra@d}%
  \pgfmathsetmacro{\xld}{\xl+\lra@d}%
  \pgfmathsetmacro{\ym}{\lra@h/2}%
  \pgfmathsetmacro{\cx}{\xl+\lra@w/2+\lra@d/2}%
  \pgfmathsetmacro{\tw}{\lra@w-\lra@d-4}%
  \ifnum#1=1 \def\lc{Series}\let\ld\lra@series\let\lico\lraicon@series\def\lkw{Series}\fi
  \ifnum#1=2 \def\lc{Volume}\let\ld\lra@volume\let\lico\lraicon@volume\def\lkw{Volume}\fi
  \ifnum#1=3 \def\lc{Book}\let\ld\lra@book\let\lico\lraicon@book\def\lkw{Book}\fi
  \ifnum#1=4 \def\lc{Chapter}\let\ld\lra@chapter\let\lico\lraicon@chapter\def\lkw{Chapter}\fi
  \ifnum#1=5 \def\lc{Section}\let\ld\lra@section\let\lico\lraicon@section\def\lkw{Section}\fi
  \ifnum#1=6 \def\lc{Topic}\let\ld\lra@topic\let\lico\lraicon@topic\def\lkw{Topic}\fi
  \edef\lfull{lra\lc Full}%
  \edef\ltint{lra\lc Tint}%
  \ifnum#1=1
    \edef\lra@poly{(\xl,0)--(\xr,0)--(\xrt,\ym)--(\xr,\lra@h)--(\xl,\lra@h)--cycle}%
  \else
    \edef\lra@poly{(\xl,0)--(\xr,0)--(\xrt,\ym)--(\xr,\lra@h)--(\xl,\lra@h)--(\xld,\ym)--cycle}%
  \fi
  \ifnum#1=\lra@N
    \expandafter\fill\expandafter[\lfull]\lra@poly;%
    \def\ltxt{white}%
    \let\lemph\bfseries
  \else
    \expandafter\fill\expandafter[\ltint]\lra@poly;%
    \edef\ltxt{\lfull}%
    \let\lemph\mdseries
  \fi
  \expandafter\lra@parse\expandafter{\ld}{\lkw}%
  \node[anchor=center,align=center,text=\ltxt,text width=\tw pt] at (\cx,\ym){%
    \iflra@icons\lico\\[-1.5pt]\fi
    {\scriptsize\sffamily\textls[60]{\MakeUppercase{\lra@eye}}}\\[-1pt]
    {\lemph\footnotesize\lra@ttl}};%
}

\newcommand{\LraCompactBreadcrumb}{%
  \lra@levelnum{\lra@current}%
  \begingroup\footnotesize\sffamily
  \def\sep{\;\textcolor{lraChapterFull}{\(\succ\)}\;}%
  \newcount\lra@compactprinted
  \lra@compactprinted=0
  \let\lra@lastcompacttitle\@empty
  \ifnum\lra@N>0 \lra@compactlevel{Series}{\lra@series}{lraSeriesFull}{}\fi
  \ifnum\lra@N>1 \lra@compactlevel{Volume}{\lra@volume}{lraVolumeFull}{}\fi
  \ifnum\lra@N>2 \lra@compactlevel{Book}{\lra@book}{lraBookFull}{}\fi
  \ifnum\lra@N>3 \lra@compactlevel{Chapter}{\lra@chapter}{lraChapterFull}{}\fi
  \ifnum\lra@N>4 \lra@compactlevel{Section}{\lra@section}{lraSectionFull}{}\fi
  \ifnum\lra@N>5 \lra@compactlevel{Topic}{\lra@topic}{lraTopicFull}{\bfseries}\fi
  \endgroup
}
\newcommand{\LraFooterBreadcrumb}{\LraCompactBreadcrumb}

\newcommand{\lra@compactlevel}[4]{%
  \expandafter\lra@parse\expandafter{#2}{#1}%
  \ifx\lra@ttl\lra@lastcompacttitle
  \else
    \ifnum\lra@compactprinted>0 \sep\fi
    \textcolor{#3}{{#4\lra@ttl}}%
    \advance\lra@compactprinted by 1%
    \let\lra@lastcompacttitle\lra@ttl
  \fi
}

\newcommand{\LraEnableFooterBreadcrumb}{%
  \pagestyle{fancy}%
  \fancyhf{}%
  \fancyfoot[L]{\LraFooterBreadcrumb}%
  \fancyfoot[R]{\thepage}%
  \renewcommand{\headrulewidth}{0pt}%
  \renewcommand{\footrulewidth}{0.4pt}%
  \fancypagestyle{plain}{%
    \fancyhf{}%
    \fancyfoot[L]{\LraFooterBreadcrumb}%
    \fancyfoot[R]{\thepage}%
    \renewcommand{\headrulewidth}{0pt}%
    \renewcommand{\footrulewidth}{0.4pt}%
  }%
}

\AtBeginDocument{%
  \let\lra@orig@section\section
  \RenewDocumentCommand{\section}{s o m}{%
    \IfBooleanTF{#1}{%
      \IfNoValueTF{#2}{\lra@orig@section*{#3}}{\lra@orig@section*[#2]{#3}}%
    }{%
      \lrameta{section={#3}, topic={}, current=section}%
      \IfNoValueTF{#2}{\lra@orig@section{#3}}{\lra@orig@section[#2]{#3}}%
    }%
  }%
  \let\lra@orig@subsection\subsection
  \RenewDocumentCommand{\subsection}{s o m}{%
    \IfBooleanTF{#1}{%
      \IfNoValueTF{#2}{\lra@orig@subsection*{#3}}{\lra@orig@subsection*[#2]{#3}}%
    }{%
      \lrameta{topic={#3}, current=topic}%
      \IfNoValueTF{#2}{\lra@orig@subsection{#3}}{\lra@orig@subsection[#2]{#3}}%
    }%
  }%
}

% ------------------------------------------------------------
% Legacy chapter breadcrumb compatibility
% ------------------------------------------------------------
\tcbset{
  breadcrumbstyle/.style={
    colback=breadcrumb, colframe=breadcrumbborder, coltitle=breadcrumbtitle,
    fonttitle=\bfseries\small, arc=2pt, boxrule=0.6pt,
    left=6pt, right=6pt, top=4pt, bottom=4pt,
    before skip=6pt, after skip=8pt,
  }
}

% \breadcrumb{<chapter subject>}{<prior display>}{<current display>}{<next display>}
% Empty prior/next are omitted together with their arrows. ASCII-only content.
\NewDocumentCommand{\breadcrumb}{m m m m}{%
  \begin{tcolorbox}[breadcrumbstyle, title={#1}]%
    \centering$%
      \ifblank{#2}{}{\text{#2}\;\to\;}%
      \textbf{#3}%
      \ifblank{#4}{}{\;\to\;\text{#4}}%
    $%
  \end{tcolorbox}%
}

% \stubstatus  -- appended after the breadcrumb on planned (stub) chapters.
\NewDocumentCommand{\stubstatus}{}{%
  \begin{tcolorbox}[breadcrumbstyle, colback=amber!10, colframe=amber]%
    \textbf{Status:} Planned%
  \end{tcolorbox}%
}

% ============================================================
% Simple, templated: where the development moves from -> to, and what it studies.
% \chapterroadmap{<prior topic>}{<next topic>}{<topics studied, prose or list>}
% prior/next are the registry neighbours (same source as \breadcrumb); the topic
% list is the chapter's section inventory -- so this can be emitted deterministically.
% ============================================================
\NewDocumentCommand{\chapterroadmap}{m m m}{%
  \begin{exposition}%
  This chapter moves the development from #1 to #2. To get there, it studies: #3.%
  \end{exposition}%
}

\makeatother
